\section{Decision \& Optimization Engine}
\label{sec:decision_optimization}

\subsection{From Forecast to Decision}
\label{subsec:from_forecast_to_decision}

While the forecasting engine predicts what is likely to happen across the 24--72 hour horizon, the Decision and Optimization Engine determines what operational responses should be considered.

Its primary role is to evaluate impending renewable generation conditions, detect operational imbalances, and synthesize candidate dispatch actions under physical system constraints. The foundational operational progression is:
\begin{center}
    \textbf{\color{primaryNavy}Forecast + Risk $\longrightarrow$ Possible Actions $\longrightarrow$ Action Evaluation $\longrightarrow$ Recommended Action}
\end{center}

\begin{figure}[htbp]
    \centering
    \begin{tikzpicture}[
        node distance=0.36cm,
        pipeBox/.style={rectangle, rounded corners=3pt, draw=accentTeal!80, fill=boxBg, text width=7.8cm, align=left, font=\scriptsize, inner sep=3.5pt},
        arrow/.style={-{Stealth[scale=0.8]}, thick, color=darkSlate}
    ]
        \node[pipeBox] (d1) {\textbf{1. Generation Forecast \& Risk Matrix} \hfill \color{darkSlate!70}Hourly Trajectory + Uncertainty Bands};
        \node[pipeBox, below=of d1] (d2) {\textbf{2. Risk \& Imbalance Event Detection} \hfill \color{darkSlate!70}Surplus/Shortfall Windows, Volatility Alerts};
        \node[pipeBox, below=of d2] (d3) {\textbf{3. Candidate Action Generation} \hfill \color{darkSlate!70}Storage Dispatch, Import/Export, Backup, Curtailment};
        \node[pipeBox, below=of d3] (d4) {\textbf{4. What-If Scenario Simulation} \hfill \color{darkSlate!70}Operator Parameter Shifts (Weather, Demand, SOC)};
        \node[pipeBox, below=of d4] (d5) {\textbf{5. Constraint \& Resource Evaluation} \hfill \color{darkSlate!70}Battery State, Ramp Rates, Interconnection Limits};
        \node[pipeBox, below=of d5] (d6) {\textbf{6. Multi-Criteria Action Comparison} \hfill \color{darkSlate!70}Reliability, Cost Impact, Carbon Emissions, Wastage};
        \node[pipeBox, below=of d6] (d7) {\textbf{7. Explainable Recommendation} \hfill \color{darkSlate!70}Structured Advice: Action + Reason + Impact};
        \node[pipeBox, below=of d7, draw=accentAmber!90, fill=accentAmber!10] (d8) {\textbf{8. Operator Decision \& Review} \hfill \color{darkSlate!70}\textbf{Human-in-the-Loop Approval \& Supervisory Dispatch}};

        \draw[arrow] (d1) -- (d2);
        \draw[arrow] (d2) -- (d3);
        \draw[arrow] (d3) -- (d4);
        \draw[arrow] (d4) -- (d5);
        \draw[arrow] (d5) -- (d6);
        \draw[arrow] (d6) -- (d7);
        \draw[arrow] (d7) -- (d8);
    \end{tikzpicture}
    \caption{Decision and Optimization Pipeline from forecast ingestion to human operator approval.}
    \label{fig:decision_pipeline}
\end{figure}

As shown in \cref{fig:decision_pipeline}, the engine acts as an advisory decision-support layer, translating complex multi-variable constraints into actionable control-room guidance.

\subsection{Risk \& Event Detection}
\label{subsec:risk_event_detection}

GridSense AI continuously scans forecast trajectories and prediction intervals to identify impending grid events before they manifest as operational emergencies. Key detected conditions include:
\begin{itemize}[leftmargin=1.8em, itemsep=0.2em]
    \item \textbf{Renewable Shortfall:} Periods where expected renewable generation falls below scheduled base demand, threatening supply-demand balance.
    \item \textbf{Renewable Surplus:} Hours where generation exceeds local demand and export limits, threatening network over-voltage or thermal congestion.
    \item \textbf{High Forecast Dispersion:} Intervals where prediction intervals widen significantly ($p_{90} - p_{10}$ exceeds tolerance), requiring higher spinning reserves.
    \item \textbf{Storage Opportunity Windows:} Surplus periods coinciding with low battery state of charge (SOC), ideal for economic energy absorption.
    \item \textbf{Backup Commitment Windows:} Severe deficit periods where available storage is insufficient, requiring advance commitment of thermal units.
    \item \textbf{Curtailment Risk Windows:} Unavoidable excess generation periods where storage is fully charged and export capacity is saturated.
\end{itemize}

Events are automatically scored and prioritized based on their projected severity, duration, and associated uncertainty.

\subsection{Action Options}
\label{subsec:action_options}

Depending on resource availability and system topology, the decision engine evaluates candidate actions tailored to the specific imbalance state (\cref{fig:surplus_vs_shortfall}):

\begin{figure}[htbp]
    \centering
    \begin{tikzpicture}[
        node distance=0.6cm and 0.8cm,
        condBox/.style={rectangle, rounded corners=3pt, draw=primaryNavy, fill=boxBg, text width=5.2cm, align=center, font=\footnotesize\bfseries, inner sep=5pt},
        branchBox/.style={rectangle, rounded corners=3pt, draw=borderGray, fill=white, text width=3.6cm, align=center, font=\scriptsize, inner sep=4pt},
        shortBox/.style={rectangle, rounded corners=3pt, draw=accentRose!80, fill=accentRose!10, text width=3.6cm, align=center, font=\scriptsize, inner sep=4pt},
        surpBox/.style={rectangle, rounded corners=3pt, draw=accentTeal!80, fill=accentTeal!10, text width=3.6cm, align=center, font=\scriptsize, inner sep=4pt},
        arrow/.style={-{Stealth[scale=0.85]}, thick, color=darkSlate}
    ]
        % Top Condition
        \node[condBox] (cond) {DETECTED RENEWABLE CONDITION};

        % Shortfall Branch (Left)
        \node[shortBox, below left=0.6cm and 0.2cm of cond] (shortHeader) {\textbf{SHORTFALL CONDITION}\\(Generation $<$ Demand)};
        \node[branchBox, below=of shortHeader] (s1) {\textbf{1. Storage Discharge}\\Dispatch co-located battery to bridge deficit};
        \node[branchBox, below=of s1] (s2) {\textbf{2. Energy Import}\\Procure grid/intertie power within line limits};
        \node[branchBox, below=of s2] (s3) {\textbf{3. Backup Activation}\\Commit standby thermal/hydro generation};

        \draw[arrow] (cond.south) -| (shortHeader.north);
        \draw[arrow] (shortHeader) -- (s1);
        \draw[arrow] (s1) -- (s2);
        \draw[arrow] (s2) -- (s3);

        % Surplus Branch (Right)
        \node[surpBox, below right=0.6cm and 0.2cm of cond] (surpHeader) {\textbf{SURPLUS CONDITION}\\(Generation $>$ Demand)};
        \node[branchBox, below=of surpHeader] (o1) {\textbf{1. Storage Charging}\\Absorb excess generation into BESS};
        \node[branchBox, below=of o1] (o2) {\textbf{2. Energy Export}\\Route surplus to adjacent grid regions};
        \node[branchBox, below=of o2] (o3) {\textbf{3. Managed Curtailment}\\Ramp down inverters/turbines if required};

        \draw[arrow] (cond.south) -| (surpHeader.north);
        \draw[arrow] (surpHeader) -- (o1);
        \draw[arrow] (o1) -- (o2);
        \draw[arrow] (o2) -- (o3);
    \end{tikzpicture}
    \caption{Candidate operational response pathways for renewable shortfall and surplus conditions (advisory options).}
    \label{fig:surplus_vs_shortfall}
\end{figure}

Crucially, GridSense AI does not assume all assets are perpetually available; recommendations dynamically adapt to the physical presence, current state, and operational constraints of each site.

\subsection{What-If Simulation}
\label{subsec:what_if_simulation}

A core differentiator of GridSense AI is its interactive **scenario simulation sandbox**. Rather than forcing operators to accept static recommendations, the platform allows operators to simulate operational stress cases and evaluate alternative futures:
\begin{itemize}[leftmargin=1.8em, itemsep=0.2em]
    \item \textbf{Atmospheric Perturbations:} \emph{``What if cloud cover increases by 30\% across the afternoon peak?''}
    \item \textbf{Wind Gradient Shifts:} \emph{``What if wind speeds drop below cut-in thresholds 2 hours earlier than forecast?''}
    \item \textbf{Demand Spikes:} \emph{``What if regional load exceeds scheduled baselines by 15~MW?''}
    \item \textbf{Storage Degradation / State Limitations:} \emph{``What if the battery system initial SOC is only 40\% instead of 80\%?''}
    \item \textbf{Asset Outages:} \emph{``What if standby thermal backup unit 2 is offline for unscheduled maintenance?''}
\end{itemize}

Upon adjusting scenario parameters, the engine recalculates the forecast trajectory, updates the risk matrix, and presents revised dispatch options in real time. This capability transforms GridSense AI from a passive monitoring screen into an active contingency planning sandbox.

\subsection{Action Evaluation}
\label{subsec:action_evaluation}

When multiple operational interventions are physically feasible, GridSense AI evaluates candidate actions across a multi-criteria decision framework:
\begin{enumerate}[leftmargin=1.8em, itemsep=0.25em]
    \item \textbf{Technical Feasibility \& Grid Adequacy:} Ability of the action to fully resolve the projected surplus or deficit without violating inverter limits, line ratings, or battery C-rates.
    \item \textbf{Renewable Energy Utilization:} Maximizing clean energy absorption and minimizing avoidable curtailment.
    \item \textbf{Operational Cost Impact:} Minimizing operating expenditures, including fuel costs of thermal backup, battery cycling degradation, and grid import tariffs.
    \item \textbf{Carbon Footprint:} Comparing the marginal carbon intensity of candidate dispatch sequences (e.g., battery discharge vs. fossil peaker activation).
    \item \textbf{System Reliability Margins:} Preserving adequate reserve capacity to absorb subsequent forecast volatility.
\end{enumerate}

This holistic evaluation ensures that recommendations reflect realistic operational, economic, and environmental trade-offs.

\subsection{Explainable Recommendation}
\label{subsec:explainable_recommendations}

To ensure full transparency in the control room, every recommendation generated by GridSense AI is structured according to a tripartite explainability standard:
\begin{center}
    \textbf{\color{primaryNavy}Recommended Action \quad + \quad Underlying Reason \quad + \quad Expected Operational Impact}
\end{center}

\noindent\textbf{Illustrative Control-Room Recommendation:}
\begin{quote}
    \small
    \textbf{Recommended Action:} Dispatch BESS at 18~MW from 17:30 to 19:30; maintain backup thermal generator on standby.\\
    \textbf{Operational Reason:} Solar generation is projected to drop below net demand starting at 17:30. Current BESS State of Charge (76\%) is sufficient to supply the projected 18~MW shortfall for 2 hours.\\
    \textbf{Expected Impact:} Eliminates the requirement to fire the fossil-fuel backup peaker plant, reducing fuel operating costs and avoiding peaker carbon emissions while maintaining frequency stability.
\end{quote}

By providing explicit justifications, GridSense AI eliminates ``black-box'' hesitation, allowing operators to verify and act with confidence.

\subsection{Decision Loop}
\label{subsec:decision_loop}

The closed-loop decision workflow proceeds through eight synchronized phases:
\begin{center}
    \texttt{Forecast} $\longrightarrow$ \texttt{Detect Risk} $\longrightarrow$ \texttt{Generate Actions} $\longrightarrow$ \texttt{Simulate Scenarios} $\longrightarrow$ \texttt{Evaluate Constraints} $\longrightarrow$ \texttt{Compare Actions} $\longrightarrow$ \texttt{Recommend} $\longrightarrow$ \textbf{\color{accentAmber}Operator Decision}
\end{center}

The human operator retains ultimate authority and supervisory responsibility for approving and dispatching actions through utility supervisory control and data acquisition (SCADA) systems.

\subsection{Phase 5 Conclusion}
\label{subsec:phase5_conclusion}

The Decision \& Optimization Engine forms the cognitive bridge between raw renewable forecasting and real-world grid operations.

By uniting proactive risk detection, what-if scenario simulations, multi-criteria action evaluations, and explainable recommendations, GridSense AI ensures that control-room operators are equipped not only to anticipate weather variability, but to execute coordinated, cost-effective, and low-carbon dispatch decisions.
